Generative artificial intelligence is increasingly capable of supporting engineering
activities that involve the interpretation of heterogeneous, weakly structured, or
partially formalized information. However, the productive use of probabilistic AI in
engineering requires more than the generation of plausible technical content. Engineering
results must remain grounded in relevant source information, uncertainty must be made
visible, consequential decisions require explicitly assigned authority, and generated
artifacts need to remain traceable and independently verifiable.

This paper derives a set of design principles for controlled AI-supported engineering
from the development and investigation of an information-processing architecture for
the transformation of heterogeneous legacy engineering information into SysML v2.
The resulting principles address the definition of the engineering task and information
basis, probabilistic interpretation, uncertainty handling, human authority, constrained
downstream processing, provenance, traceability, and independent verification.

The paper further organizes these principles into a structured path toward productive use.
The proposed perspective is not intended as a prescriptive implementation method.
Instead, it identifies the questions and control boundaries that need to be addressed when
probabilistic AI is introduced into engineering workflows. The resulting view positions the
transition from AI experimentation to engineering practice primarily as a systems and
process design challenge rather than a prompting problem.
